\chapter[INTRODUÇÃO]{INTRODUÇÃO}

Sistemas de controle industrial modernos gerenciam simultaneamente centenas ou milhares de malhas de controle realimentadas. O sistema criogênico do Grande Colisor de Hádrons (LHC) do CERN, por exemplo, opera aproximadamente 5.000 malhas PID para manter a temperatura dos imãs supercondutores a 1,9~K \cite{CERN_WEAPL02}. Em instalações dessa escala, a supervisão manual da qualidade de cada malha é inviável: controladores com ganhos inadequados, tempos de integração mal ajustados ou simplesmente lentos ou rápidos demais podem causar oscilações persistentes, desperdício de energia e, nos casos mais graves, acionamento de intertravamentos de segurança que interrompem a operação \cite{CERN_WEAPL02}. O diagnóstico automático de desempenho de malhas fechadas — identificar quais controladores estão mal sintonizados antes que causem falhas — é portanto um problema relevante tanto em instalações de grande porte quanto em plantas industriais de médio porte.

Uma abordagem promissora para desenvolver e validar estratégias de supervisão de malhas sem risco à planta real é o uso de \textbf{gêmeos digitais} (\textit{digital twins}): representações digitais compreensivas de um processo físico que replicam sua dinâmica com fidelidade suficiente para que políticas de controle e supervisão possam ser testadas, ajustadas e avaliadas em ambiente simulado \cite{Boschert2016}. Um gêmeo digital operacional permite ao engenheiro introduzir distúrbios, degradar controladores deliberadamente e observar a resposta do sistema — experimentos que seriam impraticáveis ou perigosos na planta real. Além disso, um gêmeo digital que evolui com o ativo ao longo do ciclo de vida operacional pode gerar históricos de alta qualidade para cálculo de métricas de desempenho de malhas e, eventualmente, suportar sistemas de assistência autônomos para manutenção preditiva \cite{Boschert2016}.

Para que o gêmeo digital sirva como plataforma de desenvolvimento de supervisores, é necessário um processo de referência suficientemente complexo e bem documentado. O \textbf{Processo Tennessee Eastman} (TEP), proposto por Downs e Vogel em 1993, tornou-se o \textit{benchmark} canônico da comunidade de controle de processos para esse fim \cite{DOWNS1993}. O TEP modela uma planta química realista com cinco unidades de processo, reações não-lineares acopladas, laço de reciclo e 20 distúrbios programados — combinação que desafia qualquer estratégia de controle e supervisão e permite avaliação quantitativa de desempenho. Sua ampla adoção na literatura garante que resultados obtidos sobre o TEP sejam comparáveis a trabalhos anteriores e futuros.

Plantas de processo contínuo como o TEP não podem ser operadas apenas por malhas de controle locais. A literatura consolidada de \textbf{plantwide control} \cite{Larsson2000, Skogestad2004} estabelece que é necessária uma camada supervisória com visão holística da planta, capaz de coordenar as malhas locais e reagir a distúrbios que transcendem os limites de cada controlador individual. Skogestad formaliza essa hierarquia em três camadas: controle regulatório (malhas PID locais), controle supervisório (coordenação e mudanças de setpoints), e otimização em tempo real (reconfiguração de objetivos operacionais). Essa estrutura hierárquica é exatamente aquela que o laboratório proposto neste trabalho pretende implementar e validar.

Paralelamente, o ecossistema de orquestração de contêineres consolidado em torno do \textbf{Kubernetes} oferece uma plataforma natural para hospedar componentes de supervisão industrial: o padrão de \textit{Operators} permite codificar lógica supervisória como controladores declarativos que observam o estado de uma aplicação e agem para manter a operação dentro de parâmetros definidos \cite{Burns2016}. A separação entre política (declarada em um recurso customizado) e execução (implementada no \textit{controller loop}) é análoga à separação entre especificação e runtime em arquiteturas de controle distribuído, e a infraestrutura de alta disponibilidade do Kubernetes garante que o supervisor continue operando mesmo em face de falhas de componentes individuais \cite{Burns2016}.

O presente trabalho propõe refatorar o modelo original do TEP de FORTRAN para um \textbf{gêmeo digital implementado em Rust}, servindo como plataforma de desenvolvimento e validação de estratégias de supervisão. O gêmeo digital é disponibilizado via gRPC, permitindo integração com múltiplos clientes. Uma Interface Homem-Máquina (IHM) em FastAPI com WebSocket fornece monitoramento em tempo real da planta simulada. A validação é realizada por meio de \textbf{17 experimentos} que caracterizam o comportamento da planta e testam a resposta do sistema a distúrbios programados — sendo os três últimos fundamentais para demonstração das capacidades supervisórias.

Além do digital twin validado, o trabalho \textbf{pavimenta o caminho para evolução futura} por meio de três propostas de extensão, para as quais código básico é desenvolvido mas não validado:
\begin{enumerate}
  \item \textbf{Integração com Kubernetes:} um Operator básico em Go com Kubebuilder codifica lógica supervisória declarativa, demonstrando viabilidade arquitetural sem implementação completa de estratégias de controle;
  \item \textbf{Controle \textit{plant-wide}:} proposta de hierarquia de controle (regulatório, supervisório, otimização) estruturada segundo Skogestad, com infraestrutura preparada mas não testada;
  \item \textbf{Integração OPC UA:} mapeamento do espaço de endereçamento e exposição das variáveis de processo via protocolo IEC~62541, viabilizando integração futura com sistemas industriais reais;
  \item \textbf{Índice de Preditividade:} framework para cálculo da métrica proposta pelo CERN \cite{CERN_WEAPL02}, preparado para análise de qualidade de malhas em trabalhos subsequentes.
\end{enumerate}

Esta estratégia bifurcada permite entregar um laboratório funcional e imediatamente útil, enquanto estabelece direções claras para pesquisa e desenvolvimento posteriores.


\section{PROBLEMA}
\label{sec:problema}

Desenvolver estratégias de supervisão para sistemas de controle com múltiplas malhas requer um ambiente de experimentação seguro, reprodutível e suficientemente rico em dinâmica. Plantas reais não permitem a introdução deliberada de distúrbios ou a degradação controlada de controladores para fins de teste. Simuladores genéricos carecem de realismo suficiente para validar lógica supervisória antes da implantação. Há, portanto, a necessidade de uma infraestrutura de laboratório que combine um gêmeo digital de alta fidelidade com uma camada supervisória automatizada e observável.

\section{A PROPOSTA}
\label{sec:proposta}

Este trabalho propõe refatorar o modelo original do Tennessee Eastman de FORTRAN para um gêmeo digital em Rust, servindo como plataforma de desenvolvimento e validação de estratégias de supervisão em sistemas de controle. A entrega principal compreende:

\begin{enumerate}
  \item Um simulador da planta em Rust com integração numérica RK4 e comunicação via gRPC;
  \item Uma Interface Homem-Máquina (IHM) web em FastAPI com WebSocket para monitoramento em tempo real e controle de distúrbios;
  \item Validação funcional por meio de 17 experimentos que caracterizam o comportamento da planta sob distúrbios programados, sendo os três últimos fundamentais para demonstração das capacidades de supervisão.
\end{enumerate}

Além desse núcleo validado, o trabalho estabelece a infraestrutura para futuras extensões por meio de código básico (não validado) para: (i) integração com Kubernetes via Operator declarativo; (ii) propostas de controle \textit{plant-wide} segundo a hierarquia de Skogestad; (iii) integração OPC UA para compatibilidade com sistemas industriais reais; e (iv) framework de cálculo do Índice de Preditividade para avaliação de qualidade de malhas. Esta estratégia bifurcada permite entregar um laboratório funcional e imediatamente útil, enquanto pavimenta direções claras para pesquisa subsequente.

\section{OBJETIVOS}
\label{sec:objetivos}

\subsection{Objetivo Geral}

Desenvolver e validar um gêmeo digital funcional do Processo Tennessee Eastman como plataforma para desenvolvimento, teste e validação de estratégias de supervisão de controle industrial.

\subsection{Objetivos Específicos}
\begin{enumerate}
    \item Implementar o modelo matemático do TEP em Rust com integração numérica RK4 e comunicação via gRPC.
    \item Desenvolver uma Interface Homem-Máquina (IHM) web em FastAPI com WebSocket para monitoramento em tempo real e controle de distúrbios.
    \item Validar a integração dos componentes (simulador + IHM) por meio de 17 experimentos com distúrbios programados do TEP.
    \item Caracterizar o comportamento da planta e dos controladores sob os distúrbios IDV(1) e IDV(2), identificando oportunidades para supervisão e controle avançados.
    \item Demonstrar viabilidade técnica de integração com Kubernetes através de protótipo de Operator declarativo (não validado operacionalmente).
\end{enumerate}

\section{ORGANIZAÇÃO DO TRABALHO}

Este documento está organizado da seguinte forma. O Capítulo~2 apresenta os fundamentos teóricos: o Processo Tennessee Eastman, o conceito de gêmeo digital, os Operators do Kubernetes e o protocolo gRPC. O Capítulo~3 descreve o desenvolvimento da infraestrutura proposta, detalhando cada componente e as decisões de arquitetura. O Capítulo~4 apresenta os resultados dos 17 experimentos realizados e a análise do comportamento da planta sob distúrbios. O Capítulo~5 conclui o trabalho e apresenta as direções de trabalhos futuros, incluindo a incorporação de métricas de qualidade de malha à camada supervisória.
